\subsection{一个例子}

设 \(k\) 为一个域，\(n\) 为一个 \(\geqslant 0\) 的整数，且 \((e_i)\) 为 \(k^n\) 的标准基。令 \(G = GL(n, k)\)，令 \(B\) 为 \(G\) 的上三角子群，令 \(N\) 为 \(G\) 的由每行每列恰有一个非零元素的矩阵组成的子群。\(N\) 中的一个元素置换直线 \(ke_i\)；这给出一个满同态 \(N \to \mathfrak{S}_n\)，其核是对角矩阵组成的子群 \(T = B \cap N\)，并允许我们将 \(W = N/T\) 与 \(\mathfrak{S}_n\) 认同。记 \(s_j\)（\(1 \leqslant j \leqslant n - 1\)）为 \(W\) 中对应于 \(j\) 与 \(j + 1\) 的换位的元素；令 \(S\) 为 \(s_j\) 的集合。四元组 \((G, B, N, S)\) 是一个 Tits 系统。事实上：

公理 (T1) 由《代数》第 II 章第 10 节第 13 号命题 2 的推论得出。

公理 (T2) 在《代数》第 I 章第 97 页的勘误中得到证明。

公理 (T4) 是立即的。

还需验证公理 (T3)，即

\[
s _ {j} \mathrm{B} w \subset \mathrm{B} w \mathrm{B} \cup \mathrm{B} s _ {j} w \mathrm{B} \quad \text { for } 1 \leqslant j \leqslant n - 1, w \in \mathrm{W},
\]

或等价地，

\[
s _ {j} \mathrm{B} \subset \mathrm{BB} ^ {\prime} \cup \mathrm{Bs} _ {j} \mathrm{B} ^ {\prime}, \quad \text { with } \mathrm{B} ^ {\prime} = w \mathrm{B} w ^ {- 1}.
\]

令 \(G_{j}\) 为 G 的如下子群：其中的元素固定所有 \(e_{i}\)（对于 \(i \neq j, j + 1\)），并保持由 \(e_{j}\) 和 \(e_{j+1}\) 张成的平面；该群同构于 \(\mathbf{GL}(2, k)\)。容易验证 \(G_{j}B = BG_{j}\)。由于 \(s_{j} \in G_{j}\)，有 \(s_{j}B \subset BG_{j}\)，因此只需证明

\[
\mathrm{G} _ {j} \subset (\mathrm{B} \cap \mathrm{G} _ {j}) (\mathrm{B} ^ {\prime} \cap \mathrm{G} _ {j}) \cup (\mathrm{B} \cap \mathrm{G} _ {j}) s _ {j} (\mathrm{B} ^ {\prime} \cap \mathrm{G} _ {j}).
\]

将 \(G_{j}\) 与 \(\mathbf{GL}(2,k)\) 认同；此时群 \(B \cap G_{j}\) 认同于上三角子群 \(B_{2}\)（属于 \(\mathbf{GL}(2,k)\)），而群 \(B' \cap G_{j}\) 在第一种情形下认同于 \(\mathbf{B}_2\)，其中 \(w(j) < w(j + 1)\)，在其他情形下认同于下三角子群 \(\mathbf{B}_2^-\)。在第一种情形下，要证明的公式可写成

\[
\mathbf {G L} (2, k) = \mathrm{B} _ {2} \cup \mathrm{B} _ {2} s \mathrm{B} _ {2} \quad \text { where } \quad s = \left( \begin{array}{c c} 0 & 1 \\ 1 & 0 \end{array} \right);
\]

例如，这由以下事实得出：\(B_{2}\) 是 \(\mathbf{GL}(2,k)\) 作用于射影直线 \(\mathbf{P}_{1}(k)\) 时某一点的稳定子，并且在该点的补集上传递作用。在第二种情形下，要证明的公式可写成

\[
\mathbf {G L} (2, k) = \mathrm{B} _ {2} \mathrm{B} _ {2} ^ {-} \cup \mathrm{B} _ {2} s \mathrm{B} _ {2} ^ {-};
\]

由于 \(B_{2}^{-} = sB_{2}s\)，将前一个公式右乘 s 即得此结论。

